% Nouveau EGA V (cote 25): archive cote 25, site page 69 (PDF page 70).
% typescript. Draft compared: none.
% AI-transcribed by Claude Opus 5.5 (claude-opus-5-5, adaptive thinking, effort high), Message Batches API, batch ega25b (msgbatch_01A7n5YuXubUPNK1zsepi8RG), 2026-11-13; not yet checked by eye.
\scanpage{25/69}
\mspage{69}

inchangée (cf aussi \inserted{\unsure{plus bas}} commentaires à n\textsuperscript{o} 8). 2.13. remplacer fini
par dim \hand{$\leqslant$} m-1. Th. 2.14. reste valable tel quel. , 2.15. en
remplaçant fini par dim \hand{$\leqslant$} m-1. 2.16. valable t el quel\xout{x} ,dans
\xout{2x} 2.17. remplacer fini par dim \hand{$\leqslant$} m-1. 2.18. tel quel.

Pour 3.1. , on peut l'énoncer avec m quelconque en supposant
dim f(X) \hand{$\geqslant$} m+1 , mais je propose de garder l'énoncé principal dans
le cas des \hand{?}hypersurfaces, et de signaler le cas général en corollaire
ou en remarque (il s'en\xout{x} déduit aussitôt par le procédé habituel de
prendre des hyperplans génériques indépendants). A moins qu'il soit
amusant d'expliciter une version généralisée du lemme 3.1.1. ...
Pour 3.2. lire dim f(X)\hand{$\geqslant$} m+1 , dans 3.3. remplacer dim f(X$_i$) \hand{$\geqslant$} 2
par dim f(X$_i$)\hand{$\geqslant$} m+1, et dans la\xout{x} définition de \hand{\unsure{$\mathcal{C}$|$\mathfrak{C}$}} dim f(Z) = 0
par dim f(Z) \hand{$\leqslant$} m-1.

Les considérations générales \unsure{du|da} n\textsuperscript{o} 4 s'appliquent tels quels
au cas m quelconque. Il en e st de même de 4.2. , et de 4.3. en rempla-%
çant dans b) (v) et (vi) la condition dimensionelle par
dim f(X)\hand{$\geqslant$ m+1}. Changement analogue dans 4.4. b).

Le laïus 5.1. passe tel quel. Dans 5.2. il faut se rappeler
que $\phi$ devient une section $\phi^{(m)}$ de $\underline{G}_T(1)$ (où T \xout{est un préschéma}
\xout{au dessus de} = $X\mathrm{x}_S Gr^m$) \xout{\unsure{et la condition}} induisant des sections
$\phi^m_{\hand{$\xi$}}$ des \xout{\unsure{(m)}} \xout{\illegible{1}} $\underline{O}_{X_{k(\hand{$\xi$})}}(1)\hand{$\otimes$}_{k(\hand{$\xi$})}\underline{G}(\hand{$\xi$})$ (pour \hand{$\xi \in$} Gr$^m$). Or
de façon générale, on aura expliqué au par.19 , lorsqu'on a une
\margin{$\|$} section $\phi$ d'un Module loc libre de rang m sur un préschéma, ce que
\margin{$\|$} cela veut dire que cette s ection soit $\underline{F}$ -régulière, pour un Module
\margin{$\|$} $\underline{F}$ donné : en termes d'une base locale, cela signifie qu'on a une
\margin{$\|$} suite F-régulière de m sections de $\underline{O}_X$ (et il faut vérifier que c ela

\mspage{69}
